\section{What this means for your idea}\label{sec:meaning}

Your idea has a starting claim and a prediction. The experiment bears on them differently.

\textbf{The starting claim is supported in these models.} A word's row holds most of what
the word means: section \ref{sec:probe} could read 70 to 97 in 100 of the animals' facts from
their rows alone, and animals that mean the same were given nearly the same row.

\textbf{The prediction was neither confirmed nor refuted, but there is a lead.} The
prediction is that a model of a given size does worse once it is given more words than fit.
Our rule looks for a model that comes to a stop at a lower score on more words while a bigger
model does better.
\begin{itemize}
\item With 8 bits per animal, no model ended up worse, even with 64 times as many animals. But
      by the rough guide no model was more than a little over one-third full, so the
      prediction was never put to the test.
\item With 32 bits per animal, widths 16 and 32 learned the 256 animals as well, though some
      of their runs first stood still for a long time. Width 8 did not learn them, in three
      runs on two repeats. That is the pattern of your prediction. It falls short of our rule
      because none of those runs ended flat by our test, and I cannot yet tell a model that is
      too small from one that is very slow.
\item Width 4 learned no facts on any language, so it says nothing about more words.
\end{itemize}

\textbf{Why adding words did not fill a model.} One part of the answer follows from how the
model is built: every animal comes with a row of its own. In the first round each added animal
brought 8 bits to be kept and, at width 8, 8 new numbers to keep them in: 1 bit per number,
half the rough limit, however many animals there are.

The other part is my guess. The number of words alone does not decide how big a model must be.
Two other things do. One is how much each word means, because that has to fit into the word's
own row, and a row has a fixed width. The other is how much there is to know about words
\emph{in combination} (here, an animal together with a season), because that is worked out in
the body, which all the words share. The 32-bit round speaks for this: width 8, with the
smallest rows of the three widths in that round, is the only one that did not learn the 256
animals. At width 8 a row has 8 numbers, which by the rough guide hold about 16 bits. Each
animal now has 32, so about half of what there is to know about it must be kept in the body.
For 16 animals that overflow takes about 5\% of the body; for 256 animals about four-fifths of
it, which by so rough a guide may be all the body can take. If that is right,
it is your idea in another form: the part of the model that all words share fills up as words
are added. It is a guess, and the next experiment should test it.

\textbf{What a larger vocabulary did cost: training.} With 8 bits per animal, 64 times as many
animals took about 6 to 8 times as many sets. With 32 bits per animal, 16 times as many animals
took 4 to 6 times as many sets when a run went straight through (widths 16 and 32: 91 and 64
sets against 16), and more than 20 times as many when it stood still (more than 362).

\textbf{A lesson about method.} During a stand-still a model looks exactly like one that has
reached its limit, and our stopping test was fooled by that on two occasions, the second
after I had already repaired it once. Your later suggestion, to give every model the same amount of text
and then compare, needs no judgment about when a model has finished. With amounts chosen
generously, I would use it in the next experiment.

\textbf{What this experiment cannot tell you.}
\begin{itemize}
\item It is a toy: the ``meaning'' of a word is four or sixteen facts, and the models are
      many thousands of times smaller than real ones. Nothing here carries over to large
      language models without further work.
\item The score is memory of facts. It says nothing about fluency or reasoning.
\item No score below 100 in this report is a proven limit. ``Flat'' and ``not finished'' both
      mean only that the model had not got further when its run ended.
\item Much of the 32-bit round rests on one or two runs for each setting, and that round
      showed how much one run can differ from the next.
\end{itemize}
